problem:  
Let \( X = G_1/K_1 \) and \( Y = G_2/K_2 \) be irreducible symmetric spaces of noncompact type with \( \operatorname{rank}(X), \operatorname{rank}(Y) \ge 2 \).  
Let \( f: X \to Y \) be a smooth harmonic map satisfying the bounded energy density condition  
\[
\sup_{x \in X} \|df(x)\| < \infty .
\]  
For each \(x \in X\), denote by  
\[
\mathrm{rank}_X(f;x) = \sup_{F} \dim(df_x(T_xF)),
\]  
where the supremum is taken over all **maximal flats** \(F \subset X\) passing through \(x\).  

**Problem:**  
Must every non–totally geodesic harmonic map \(f: X \to Y\) satisfy  
\[
\sup_{x \in X}\mathrm{rank}_X(f;x) = \operatorname{rank}(Y) \,?
\]
Equivalently, can a harmonic map between higher-rank symmetric spaces remain **rank-deficient along all flats**, i.e. have  
\(\sup_{x}\mathrm{rank}_X(f;x) < \operatorname{rank}(Y)\),  
without being totally geodesic?  

Moreover, if the equality always holds, does this phenomenon admit an explanation through a **Bochner-type curvature identity** on \(X\) and \(Y\) that analytically enforces full-rank behavior along some flat direction?  

---

Why is it a "good" problem:  
This problem isolates a precise and intrinsic **local rank-rigidity question** for harmonic maps between higher-rank symmetric spaces. Unlike global superrigidity or totally geodesic results that require group actions or topological constraints, this asks whether **harmonicity itself**—as a differential analytic condition—prevents infinitesimal rank collapse along every Cartan direction.  

It sits at the intersection of  
- **geometric analysis** (Bochner identities, curvature terms controlling harmonic maps),  
- **differential geometry of symmetric spaces** (maximal flats, rank structures), and  
- **Lie theory** (Cartan subspaces and root-space decompositions).  

A positive answer would reveal a new kind of *analytic rank rigidity*, showing that harmonic maps intrinsically detect the target’s algebraic rank through curvature-based differential inequalities. A counterexample would expose unexpected analytic flexibility in the higher-rank setting.  

The formulation is concise and geometric—expressed in terms of pointwise differential rank—but touches deep interactions between curvature, PDEs, and Lie symmetries. It exemplifies the ideal of a “narrow entrance, profound depth” problem: easy to state, yet potentially leading to a richer understanding of how harmonicity encodes algebraic and geometric rank.